% ============================================================================
% Model_Equations.tex
%
% Complete equation catalogue for the MHI hierarchical hydrogen permeation
% model. Section A documents Model A (bulk transport hierarchy, levels
% L1-L5, driven by Application/Proposal.ipynb). Section B documents Model B
% (surface-kinetics extension, level L6 combined with the bulk levels,
% driven by Application/Surface_proposal.ipynb).
%
% Every equation is transcribed from the Python calculation modules in
% calculations/ (the code is the source of truth). Code locations are given
% as file -> function. Compile with pdflatex + bibtex (or upload the folder
% to Overleaf).
% ============================================================================
\documentclass[11pt,a4paper]{article}

\usepackage[margin=2.2cm]{geometry}
\usepackage{amsmath,amssymb}
\usepackage{booktabs}
\usepackage{longtable}
\usepackage{array}
\usepackage[round]{natbib}
\usepackage{xcolor}
\usepackage[colorlinks=true,linkcolor=blue!60!black,citecolor=blue!60!black,urlcolor=blue!60!black]{hyperref}

% ---------------------------------------------------------------- notation macros
\newcommand{\Pup}{P_{\mathrm{up}}}
\newcommand{\Pdown}{P_{\mathrm{down}}}
\newcommand{\Pint}{P_{\mathrm{int}}}
\newcommand{\Ks}{K_{\mathrm{s}}}
\newcommand{\Kox}{K_{\mathrm{ox}}}
\newcommand{\Keq}{K_{\mathrm{eq}}}
\newcommand{\Dm}{D_{\mathrm{m}}}
\newcommand{\Dox}{D_{\mathrm{ox}}}
\newcommand{\Deff}{D_{\mathrm{eff}}}
\newcommand{\Lm}{L_{\mathrm{m}}}
\newcommand{\Lox}{L_{\mathrm{ox}}}
\newcommand{\kdiss}{k_{\mathrm{diss}}}
\newcommand{\krec}{k_{\mathrm{recomb}}}
\newcommand{\fgb}{f_{\mathrm{gb}}}
\newcommand{\Rsurf}{R_{\mathrm{surface}}}
\newcommand{\Rox}{R_{\mathrm{oxide}}}
\newcommand{\Rmet}{R_{\mathrm{metal}}}
% code identifiers (file and function names)
\newcommand{\code}[1]{\texttt{#1}}

\title{Equations of the Hierarchical Hydrogen Permeation Model\\
\large Model A: Bulk Transport Hierarchy (L1--L5) \quad
Model B: Surface-Kinetics Extension (L6)}
\author{Azeez Akinyemi}
\date{\today}

\begin{document}
\maketitle
\tableofcontents
\clearpage

% ============================================================================
\section{Introduction and notation}
% ============================================================================

The permeation model is built as a hierarchy of physical fidelity. \textbf{Model
A} (Section~\ref{sec:modelA}, notebook \code{Application/Proposal.ipynb})
assumes that surface processes are infinitely fast, so transport is always
limited by bulk diffusion through the oxide and/or metal layers. It progresses
from a perfect metal membrane (L1) through a perfect oxide (L2a), the coupled
oxide--metal bilayer (L2b), a defective oxide with parallel leakage paths (L3),
a defective metal with grain-boundary transport and trapping (L4), up to the
full defective system (L5). \textbf{Model B} (Section~\ref{sec:modelB},
notebook \code{Application/Surface\_proposal.ipynb}) relaxes the fast-surface
assumption: a finite Langmuir--Hinshelwood dissociation/recombination step (L6)
is placed in series with each bulk configuration, and the steady state is found
by enforcing flux continuity through all layers.

All equations below are transcribed directly from the Python modules in
\code{calculations/}; the code location column in the summary tables gives
\emph{file $\rightarrow$ function}. Table~\ref{tab:notation} collects the
global symbols. Throughout, $R = 8.314\;\mathrm{J\,mol^{-1}\,K^{-1}}$ is the
gas constant and $T$ the absolute temperature in K.

\begin{table}[h]
\centering
\caption{Global notation and units.}
\label{tab:notation}
\begin{tabular}{llll}
\toprule
Symbol & Meaning & Units & First used \\
\midrule
$J$ & permeation flux & $\mathrm{mol\,m^{-2}\,s^{-1}}$ & Eq.~\eqref{eq:fick} \\
$C$ & dissolved hydrogen concentration & $\mathrm{mol\,m^{-3}}$ & Eq.~\eqref{eq:sieverts} \\
$\Pup,\,\Pdown$ & upstream / downstream $\mathrm{H_2}$ pressure & Pa & Eq.~\eqref{eq:sieverts} \\
$\Pint$ & oxide--metal interface pressure & Pa & Eq.~\eqref{eq:closureL2b} \\
$\Dm,\,\Dox$ & diffusivity of H in metal / $\mathrm{H_2}$ in oxide & $\mathrm{m^2\,s^{-1}}$ & Eq.~\eqref{eq:fick} \\
$\Ks$ & Sieverts' solubility constant (metal) & $\mathrm{mol\,m^{-3}\,Pa^{-1/2}}$ & Eq.~\eqref{eq:sieverts} \\
$\Kox$ & oxide solubility constant (Henry form in A, & $\mathrm{mol\,m^{-3}\,Pa^{-1}}$ (A) & Eq.~\eqref{eq:henry} \\
 & \quad Sieverts form in B; \S\ref{sec:oxideconvention}) & $\mathrm{mol\,m^{-3}\,Pa^{-1/2}}$ (B) & \\
$\Lm,\,\Lox$ & metal / oxide thickness & m & Eq.~\eqref{eq:L1flux} \\
$\Phi$ & metal permeability $\Dm\Ks$ & $\mathrm{mol\,m^{-1}\,s^{-1}\,Pa^{-1/2}}$ & Eq.~\eqref{eq:permeability} \\
$\Deff$ & effective (microstructure-modified) diffusivity & $\mathrm{m^2\,s^{-1}}$ & Eq.~\eqref{eq:orianiDeff} \\
$\theta$ & surface coverage (fraction of occupied sites) & -- & Eq.~\eqref{eq:ads} \\
$\kdiss$ & dissociative adsorption rate constant & $\mathrm{mol\,m^{-2}\,s^{-1}\,Pa^{-1}}$ & Eq.~\eqref{eq:ads} \\
$\Keq$ & surface equilibrium constant $\kdiss/\krec$ & $\mathrm{Pa^{-1}}$ & Eq.~\eqref{eq:des} \\
$\alpha$ & oxide permeance $\Dox\Kox/\Lox$ (Model B) & $\mathrm{mol\,m^{-2}\,s^{-1}\,Pa^{-1/2}}$ & Eq.~\eqref{eq:permeances} \\
$\beta$ & metal permeance $\Dm\Ks/\Lm$ & $\mathrm{mol\,m^{-2}\,s^{-1}\,Pa^{-1/2}}$ & Eq.~\eqref{eq:permeances} \\
$E_D,\,H_s,\,E$ & activation / solution energies & $\mathrm{J\,mol^{-1}}$ & Eq.~\eqref{eq:arrhenius} \\
\bottomrule
\end{tabular}
\end{table}

% ============================================================================
\section{Model A: bulk transport hierarchy (L1--L5)}
\label{sec:modelA}
% ============================================================================

% ----------------------------------------------------------------------------
\subsection{Summary table}
% ----------------------------------------------------------------------------

{\footnotesize
\setlength{\tabcolsep}{4pt}
\begin{longtable}{@{}c >{\raggedright\arraybackslash}p{2.1cm} l >{\raggedright\arraybackslash}p{4.6cm} >{\raggedright\arraybackslash}p{2.5cm}@{}}
\caption{Model A equations, code locations (file, function in \code{calculations/}), and references.}
\label{tab:summaryA}\\
\toprule
Eq. & Name & Formula & Code location & Ref. \\
\midrule
\endfirsthead
\toprule
Eq. & Name & Formula & Code location & Ref. \\
\midrule
\endhead
\bottomrule
\endfoot
\multicolumn{5}{@{}l}{\emph{L1 --- perfect metal}}\\
\eqref{eq:sieverts} & Sieverts' law & $C = \Ks\sqrt{P}$ & \code{permeation\_calc.py}\newline \code{sieverts\_concentration} & \citet{sieverts1929} \\
\eqref{eq:fick} & Fick's first law & $J = D\,(C_{\mathrm{up}}-C_{\mathrm{down}})/L$ & \code{permeation\_calc.py}\newline \code{fick\_flux} & \citet{fick1855} \\
\eqref{eq:L1flux} & L1 flux & $J_{\mathrm{L1}} = \frac{\Dm\Ks}{\Lm}(\sqrt{\Pup}-\sqrt{\Pdown})$ & \code{permeation\_calc.py}\newline \code{calculate\_simple\_metal\_flux} & \\
\eqref{eq:permeability} & Permeability & $\Phi = \Dm\Ks$ & \code{utils.py}\newline \code{get\_permeability} & \\
\midrule
\multicolumn{5}{@{}l}{\emph{L2a --- perfect oxide}}\\
\eqref{eq:henry} & Henry's law & $C = \Kox P$ & \code{oxide\_permeation.py}\newline \code{molecular\_diffusion\_flux} & \\
\eqref{eq:L2aflux} & L2a flux & $J_{\mathrm{L2a}} = \frac{\Dox\Kox}{\Lox}(\Pup-\Pdown)$ & \code{oxide\_permeation.py}\newline \code{molecular\_diffusion\_flux} & \\
\eqref{eq:Rox} & Oxide resistance & $\Rox = \Lox/(\Dox\Kox)$ & \code{oxide\_permeation.py}\newline \code{calculate\_oxide\_resistance} & \\
\midrule
\multicolumn{5}{@{}l}{\emph{L2b --- oxide + metal bilayer}}\\
\eqref{eq:closureL2b} & Flux continuity & $J_{\mathrm{ox}}(\Pint) = J_{\mathrm{m}}(\Pint)$ & \code{interface\_solver.py}\newline \code{flux\_balance\_equation} & \\
\eqref{eq:Rmetal} & Linearised metal resistance & $\Rmet \approx 2\Lm\sqrt{\Pint}/(\Dm\Ks)$ & \code{oxide\_permeation.py}\newline \code{calculate\_metal\_resistance} & \\
\eqref{eq:chi} & Regime ratio & $\chi = \Rox/\Rmet$ & \code{interface\_solver.py}\newline \code{calculate\_oxide\_metal\_system} & \\
\eqref{eq:Ptrans} & Transition pressure & $P_{\mathrm{trans}} = \left[\frac{\Lox \Dm\Ks}{2\Dox\Kox\Lm}\right]^2$ & \code{oxide\_permeation.py}\newline \code{calculate\_transition\_pressure} & \\
\midrule
\multicolumn{5}{@{}l}{\emph{L3 --- defective oxide (parallel paths)}}\\
\eqref{eq:L3parallel} & Parallel-path flux & $J_{\mathrm{L3}} = f_{\mathrm{int}}J_{\mathrm{int}} + f_{\mathrm{def}}J_{\mathrm{def}}$ & \code{parallel\_oxide\_defect\_paths.py}\newline \code{calculate\_parallel\_path\_flux} & \citet{strehlow1974,zarchy1979} \\
\eqref{eq:crack} & Crack path & $\Lox^{\mathrm{crack}} = \gamma\,\Lox$ & \code{parallel\_oxide\_defect\_paths.py}\newline \code{calculate\_defect\_path\_flux} & \citet{strehlow1974} \\
\eqref{eq:gbdefect} & Oxide-GB path & $\Dox^{\mathrm{gb}} = \delta_D\,\Dox$ & \code{parallel\_oxide\_defect\_paths.py}\newline \code{calculate\_defect\_path\_flux} & \\
\eqref{eq:mixed} & Mixed defects & $J_{\mathrm{def}} = \sum_i w_i J_i$ & \code{parallel\_oxide\_defect\_paths.py}\newline \code{calculate\_defect\_path\_flux} & \\
\midrule
\multicolumn{5}{@{}l}{\emph{L4 --- defective metal (grain boundaries + traps)}}\\
\eqref{eq:fgb} & GB volume fraction & $\fgb = 3\delta/d$ & \code{defective\_metal.py}\newline \code{grain\_boundary\_density} & \citet{underwood1970,palumbo1990} \\
\eqref{eq:gbenh} & GB enhancement & $\alpha_{\mathrm{gb}}(T) = A\exp\!\left[\frac{Q_{\mathrm{bulk}}-Q_{\mathrm{gb}}}{RT}\right]$ & \code{defective\_metal.py}\newline \code{gb\_enhancement\_factor} & \citet{tsuru1982,brass1996} \\
\eqref{eq:gbparallel} & GB parallel model & $D_{\mathrm{gb\text{-}enh}} = (1-\fgb)D + \fgb\,\alpha_{\mathrm{gb}}D$ & \code{defective\_metal.py}\newline \code{calculate\_gb\_enhanced\_}\newline\code{diffusivity} & \\
\eqref{eq:hart} & Hart equation & see text & \code{defective\_metal.py}\newline \code{calculate\_gb\_enhanced\_}\newline\code{diffusivity} & \citet{hart1957} \\
\eqref{eq:orianitheta} & Trap occupancy & $\theta_T = \frac{K\theta_L}{1+K\theta_L}$ & \code{defective\_metal.py}\newline \code{trap\_occupancy} & \citet{oriani1970,mcnabb1963} \\
\eqref{eq:orianiDeff} & Trapping reduction & $\Deff = \frac{D}{1+\sum_i N_{T,i}K_i/N_L}$ & \code{defective\_metal.py}\newline \code{calculate\_effective\_}\newline\code{diffusivity\_trapping} & \citet{oriani1970} \\
\eqref{eq:combined} & Combined model & $\Deff = \frac{(1-\fgb)D+\fgb\alpha_{\mathrm{gb}}D}{1+\sum_i N_{T,i}K_i/N_L}$ & \code{defective\_metal.py}\newline \code{combined\_microstructure\_}\newline\code{model} & \citet{oudriss2012} \\
\eqref{eq:vacancy} & Vacancy concentration & $C_v = N\exp(-E_f^v/RT)$ & \code{defective\_metal.py}\newline \code{vacancy\_concentration} & \citet{kraftmakher1998} \\
\eqref{eq:L4flux} & L4 flux & $J_{\mathrm{L4}} = \frac{\Deff\Ks}{\Lm}(\sqrt{\Pup}-\sqrt{\Pdown})$ & \code{permeation\_calc.py}\newline \code{calculate\_defective\_}\newline\code{metal\_flux} & \\
\midrule
\multicolumn{5}{@{}l}{\emph{L5 --- full system}}\\
\eqref{eq:closureL5} & Coupled closure & $J_{\mathrm{ox}}(\Pint) = J_{\mathrm{L4}}(\Pint;\Deff)$ & \code{interface\_solver.py}\newline \code{solve\_interface\_pressure\_}\newline\code{defective\_metal} & \\
\midrule
\multicolumn{5}{@{}l}{\emph{Shared temperature dependence}}\\
\eqref{eq:arrhenius} & Arrhenius (reference form) & $k(T) = k_{\mathrm{ref}}\exp\!\left[-\frac{E}{R}\!\left(\frac1T-\frac1{T_{\mathrm{ref}}}\right)\right]$ & \code{utils.py}\newline \code{arrhenius} & \\
\eqref{eq:QPhi} & Permeation activation & $Q_\Phi = E_D + H_s$ & \code{utils.py}\newline \code{get\_permeability} & \\
\end{longtable}
}

% ----------------------------------------------------------------------------
\subsection{L1: perfect metal}
% ----------------------------------------------------------------------------

\paragraph{Sieverts' law.} At a clean metal surface in equilibrium with
$\mathrm{H_2}$ gas, the molecule dissociates and dissolves atomically, so the
dissolved concentration scales with the \emph{square root} of pressure:
\begin{equation}
  C = \Ks \sqrt{P}.
  \label{eq:sieverts}
\end{equation}
Here $C$ is the hydrogen concentration at the surface
[$\mathrm{mol\,m^{-3}}$], $P$ the hydrogen partial pressure [Pa], and $\Ks$
the Sieverts solubility constant [$\mathrm{mol\,m^{-3}\,Pa^{-1/2}}$], a
material property with its own van~'t Hoff temperature dependence
(Eq.~\eqref{eq:arrhenius} with $E = H_s$). The square root is the fingerprint
of dissociative absorption \citep{sieverts1929}: it is what distinguishes
metals from oxides in every log--log flux plot in the model.

\paragraph{Fick's first law.} Steady-state diffusion through a flat membrane
with a linear concentration profile gives
\begin{equation}
  J = D\,\frac{C_{\mathrm{up}} - C_{\mathrm{down}}}{L},
  \label{eq:fick}
\end{equation}
where $J$ is the flux [$\mathrm{mol\,m^{-2}\,s^{-1}}$], $D$ the diffusion
coefficient [$\mathrm{m^2\,s^{-1}}$], $L$ the membrane thickness [m], and
$C_{\mathrm{up}}$, $C_{\mathrm{down}}$ the boundary concentrations
[$\mathrm{mol\,m^{-3}}$]. Positive flux runs from upstream to downstream
\citep{fick1855}.

\paragraph{L1 permeation flux.} Substituting Eq.~\eqref{eq:sieverts} at both
faces into Eq.~\eqref{eq:fick} yields the classical Richardson form,
\begin{equation}
  J_{\mathrm{L1}} = \frac{\Dm \Ks}{\Lm}\left(\sqrt{\Pup} - \sqrt{\Pdown}\right),
  \label{eq:L1flux}
\end{equation}
with $\Dm$ the metal diffusivity [$\mathrm{m^2\,s^{-1}}$], $\Lm$ the metal
thickness [m], and $\Pup$, $\Pdown$ the boundary pressures [Pa]
\citep{richardson1904,sieverts1929}. On a log--log
plot of $J$ versus $\Pup$ (with $\Pdown \approx 0$) this gives the
characteristic slope of $0.5$; the notebooks use this slope as the L1
validation criterion.

\paragraph{Permeability.} The kinetic and thermodynamic factors are lumped into
\begin{equation}
  \Phi = \Dm \Ks,
  \label{eq:permeability}
\end{equation}
with units $\mathrm{mol\,m^{-1}\,s^{-1}\,Pa^{-1/2}}$. Written with $\Phi$,
Eq.~\eqref{eq:L1flux} becomes $J = (\Phi/\Lm)\,\Delta\sqrt{P}$, and the
corresponding driving-force resistance is $R_{\mathrm{L1}} = \Lm/\Phi$
[$\mathrm{Pa^{1/2}\,s\,m^{2}\,mol^{-1}}$]. Note this resistance acts on a
$\sqrt{P}$ driving force, so it is \emph{not} directly comparable to the
oxide resistance of Eq.~\eqref{eq:Rox} without linearisation
(Eq.~\eqref{eq:Rmetal}).

% ----------------------------------------------------------------------------
\subsection{L2a: perfect oxide}
% ----------------------------------------------------------------------------

\paragraph{Henry's law.} In the oxide, hydrogen is assumed to diffuse as
intact $\mathrm{H_2}$ molecules (no dissociation), so solubility is
\emph{linear} in pressure:
\begin{equation}
  C = \Kox P,
  \label{eq:henry}
\end{equation}
where $\Kox$ is the Henry constant [$\mathrm{mol\,m^{-3}\,Pa^{-1}}$] --- note
the different units from $\Ks$. This single change of exponent propagates
through the whole hierarchy: it is why oxide-limited and metal-limited regimes
are distinguishable by pressure slope.

\paragraph{L2a flux.} Combining Eq.~\eqref{eq:henry} with Fick's law across
the oxide thickness $\Lox$:
\begin{equation}
  J_{\mathrm{L2a}} = \frac{\Dox \Kox}{\Lox}\left(\Pup - \Pdown\right),
  \label{eq:L2aflux}
\end{equation}
with $\Dox$ the molecular diffusivity of $\mathrm{H_2}$ in the oxide
[$\mathrm{m^2\,s^{-1}}$] \citep{strehlow1974,hollenberg1995}. The flux is
linear in $\Delta P$ (log--log slope $1.0$), the L2a validation criterion.
The $\mathrm{Cr_2O_3}$ transport parameters themselves come from
\citet{nemanic2023} ($\Phi_{\mathrm{ox}}$, $\Dox$), \citet{stover1986}
($Q_{p,\mathrm{ox}}$), and \citet{chen2011} ($E_{D,\mathrm{ox}}$ from DFT).

\paragraph{Oxide resistance.} Because transport is linear, a genuine
pressure-independent resistance exists:
\begin{equation}
  \Rox = \frac{\Lox}{\Dox \Kox},
  \qquad J = \frac{\Delta P}{\Rox},
  \label{eq:Rox}
\end{equation}
with $\Rox$ in $\mathrm{Pa\,s\,m^{2}\,mol^{-1}}$. This constancy is the key
structural difference from the metal layer and is what makes the series
combination in L2b a nonlinear problem.

% ----------------------------------------------------------------------------
\subsection{L2b: oxide + metal in series}
% ----------------------------------------------------------------------------

\paragraph{Flux continuity closure.} With an oxide on the upstream face of the
metal, an unknown interface pressure $\Pint$ develops at the oxide--metal
boundary. At steady state the same flux must cross both layers:
\begin{equation}
  \underbrace{\frac{\Dox\Kox}{\Lox}\bigl(\Pup - \Pint\bigr)}_{J_{\mathrm{ox}}(\Pint)}
  \;=\;
  \underbrace{\frac{\Dm\Ks}{\Lm}\bigl(\sqrt{\Pint} - \sqrt{\Pdown}\bigr)}_{J_{\mathrm{m}}(\Pint)} .
  \label{eq:closureL2b}
\end{equation}
This series-barrier construction is the standard treatment of oxide permeation
barriers on structural metals \citep{strehlow1974,hollenberg1995}. The mixed
linear/square-root equation has no closed-form solution; the code
brackets $\Pint \in (\Pdown, \Pup)$ and solves it with Brent's method
(\code{interface\_solver.py} $\rightarrow$ \code{solve\_interface\_pressure}).
The solved $\Pint$ determines both the flux and the operating regime.

\paragraph{Linearised metal resistance.} To compare the two layers on a common
$\Delta P$ basis, the metal flux is linearised about $\Pint$ using
$\partial\sqrt{P}/\partial P = 1/(2\sqrt{P})$:
\begin{equation}
  \Rmet \approx \frac{2 \Lm \sqrt{\Pint}}{\Dm \Ks},
  \label{eq:Rmetal}
\end{equation}
in $\mathrm{Pa\,s\,m^{2}\,mol^{-1}}$. Unlike $\Rox$, this resistance grows
with $\sqrt{\Pint}$: at high pressure the metal becomes comparatively easier
to traverse per pascal of driving force.

\paragraph{Regime classification.} The resistance ratio
\begin{equation}
  \chi = \frac{\Rox}{\Rmet}
  \qquad
  \begin{cases}
    \chi > 10 & \text{oxide-limited } (\Pint \ll \Pup,\ \text{slope} \to 1.0),\\[2pt]
    0.5 \le \chi \le 10 & \text{transition } (0.5 < \text{slope} < 1.0),\\[2pt]
    \chi < 0.5 & \text{metal-limited } (\Pint \to \Pup,\ \text{slope} \to 0.5),
  \end{cases}
  \label{eq:chi}
\end{equation}
is evaluated at the solved interface pressure. Setting $\Rox = \Rmet$ and
solving for pressure gives the approximate crossover
\begin{equation}
  P_{\mathrm{trans}} = \left(\frac{\Lox\, \Dm \Ks}{2\, \Dox \Kox\, \Lm}\right)^{2},
  \label{eq:Ptrans}
\end{equation}
the pressure at which the rate-limiting layer switches --- a useful design
number for choosing operating windows.

\paragraph{Remark (comparing metal and oxide resistances --- the chemistry
behind the units).} The mismatch between the two resistances is not a
mathematical artefact; it records \emph{which chemical species each layer
transports}. The oxide carries \emph{intact} $\mathrm{H_2}$: no bond is
broken on dissolution, one gas molecule becomes one dissolved molecule, so
$C \propto P$ (Henry) and the flux responds linearly to $\Delta P$ with the
constant resistance $\Rox$ [$\mathrm{Pa\,s\,m^2\,mol^{-1}}$]. Dissociation is
deferred to the oxide--metal interface, where
$\mathrm{H_2} \rightleftharpoons 2\mathrm{H}^{*}$; the law of mass action for
a reaction delivering \emph{two} atoms per molecule gives
$C_{\mathrm{H}}^{2} \propto \Pint$, i.e.\ $C_{\mathrm{H}} \propto
\sqrt{\Pint}$ --- Sieverts' law \emph{is} that 2:1 stoichiometry. The metal
therefore responds to $\Delta\sqrt{P}$ (natural resistance
$R_{\mathrm{L1}} = \Lm/\Phi$ [$\mathrm{Pa^{1/2}\,s\,m^2\,mol^{-1}}$]):
doubling the pressure less than doubles the dissolved-atom population, so per
pascal the metal becomes progressively harder to feed while the oxide does
not. Eq.~\eqref{eq:Rmetal} translates the metal onto the oxide's per-pascal
basis by linearising about $\Pint$ (the factor 2 is
$\mathrm{d}\sqrt{P}/\mathrm{d}P = 1/(2\sqrt{P})$, the marginal yield of the
dissociation equilibrium), which is what makes $\chi$ in Eq.~\eqref{eq:chi}
meaningful. Two caveats: the linearised $\Rmet$ is a \emph{diagnostic only}
--- the L2b flux is always computed from the exact continuity condition
Eq.~\eqref{eq:closureL2b}, never as $\Delta P/(\Rox + \Rmet)$ --- and Model B
never meets the mismatch because there dissociation happens at the
\emph{gas--oxide} surface instead, so atomic H traverses every layer and all
driving forces are already in $\sqrt{P}$ units
(\S\ref{sec:oxideconvention}). Appendix~\ref{app:resistance} gives the
complete derivation.

% ----------------------------------------------------------------------------
\subsection{L3: defective oxide (parallel paths)}
% ----------------------------------------------------------------------------

\paragraph{Area-weighted parallel paths.} Following
\citet{strehlow1974}, a real oxide is treated as intact film and defects
acting as parallel permeation paths, weighted by their area fractions:
\begin{equation}
  J_{\mathrm{L3}} = f_{\mathrm{intact}}\, J_{\mathrm{intact}}
                  + f_{\mathrm{defect}}\, J_{\mathrm{defect}},
  \qquad f_{\mathrm{intact}} + f_{\mathrm{defect}} = 1,
  \label{eq:L3parallel}
\end{equation}
where $f_{\mathrm{defect}}$ is the fraction of surface area occupied by
defects [--] and each $J$ is a flux \emph{density} through its own path
[$\mathrm{mol\,m^{-2}\,s^{-1}}$]. Crucially, the intact oxide still conducts
($J_{\mathrm{intact}}$ from the full L2b solver), so even $\sim$1\% defect
coverage can dominate without the intact film being ignored
\citep{zarchy1979}. Each defect type maps onto a modified L1 or L2b problem:
\begin{align}
  \text{pinhole:}&\quad J_{\mathrm{def}} = J_{\mathrm{L1}}
    \quad\text{(bare metal, Eq.~\eqref{eq:L1flux}; no oxide barrier)},
    \nonumber\\
  \text{crack:}&\quad \Lox^{\mathrm{crack}} = \gamma\,\Lox,
    \quad 0<\gamma<1 \text{ (default } \gamma=0.1\text{)},
    \label{eq:crack}\\
  \text{oxide grain boundary:}&\quad \Dox^{\mathrm{gb}} = \delta_D\,\Dox,
    \quad \delta_D>1 \text{ (default } \delta_D=10\text{)},
    \label{eq:gbdefect}
\end{align}
where $\gamma$ is the crack thickness factor [--] and $\delta_D$ the
grain-boundary diffusivity enhancement [--]; crack and GB paths are solved
with the L2b machinery using the modified oxide properties. Several defect
types can coexist; their contribution is the normalised weighted sum
\begin{equation}
  J_{\mathrm{defect}} = \frac{\sum_i w_i J_i}{\sum_i w_i},
  \label{eq:mixed}
\end{equation}
with $w_i$ the component area fractions within the defective region.

% ----------------------------------------------------------------------------
\subsection{L4: defective metal (grain boundaries + trapping)}
% ----------------------------------------------------------------------------

\paragraph{Grain-boundary volume fraction.} From stereology of a random
polycrystal \citep{underwood1970}, the GB surface area per volume is
$S_v = s/d$ and the GB volume fraction is
\begin{equation}
  \fgb = \frac{s\,\delta}{d}, \qquad
  s = \begin{cases} 3 & \text{equiaxed grains}\\ 2 & \text{columnar}\\ 1 & \text{planar/lamellar}, \end{cases}
  \label{eq:fgb}
\end{equation}
where $\delta$ is the GB thickness [m] (default $0.5\;\mathrm{nm}$), $d$ the
mean grain diameter [m], and $s$ the dimensionality shape factor [--]. The
same geometry supplies the GB \emph{trap} density
$N_{\mathrm{gb}} = S_v\,\rho_{\mathrm{gb}}$ [$\mathrm{m^{-3}}$], with
$\rho_{\mathrm{gb}} \approx 10^{19}\;\mathrm{m^{-2}}$ the areal density of
trap sites on the boundary.

\paragraph{Grain-boundary enhancement factor.} GBs are fast diffusion paths
because their activation energy is lower ($Q_{\mathrm{gb}} \approx
0.6\,Q_{\mathrm{bulk}}$):
\begin{equation}
  \alpha_{\mathrm{gb}}(T) = \frac{D_{\mathrm{gb}}}{D_{\mathrm{bulk}}}
  = A \exp\!\left[\frac{Q_{\mathrm{bulk}} - Q_{\mathrm{gb}}}{RT}\right],
  \label{eq:gbenh}
\end{equation}
where $A$ is the pre-exponential ratio $D_{0,\mathrm{gb}}/D_{0,\mathrm{bulk}}$
[--] and $Q$ are activation energies [$\mathrm{J\,mol^{-1}}$]
\citep{tsuru1982,brass1996}. In the code the factor is interpolated
log-linearly in $1/T$ from tabulated austenitic-steel data
($\alpha_{\mathrm{gb}} \approx 500$ at $600\,^\circ$C down to $\approx 20$ at
$1000\,^\circ$C), then scaled by GB character (HAGB $1.0$, LAGB $0.1$, twin
$0.05$, special $0.3$).

\paragraph{GB-enhanced diffusivity.} Treating lattice and GB as independent
parallel channels:
\begin{equation}
  D_{\mathrm{gb\text{-}enh}}
  = (1-\fgb)\,D_{\mathrm{bulk}} + \fgb\,\alpha_{\mathrm{gb}} D_{\mathrm{bulk}}
  = D_{\mathrm{bulk}}\bigl[1 + \fgb(\alpha_{\mathrm{gb}}-1)\bigr].
  \label{eq:gbparallel}
\end{equation}
For larger GB fractions ($\fgb \gtrsim 0.01$), where network connectivity
matters, the Hart equation \citep{hart1957} is available instead:
\begin{equation}
  \frac{\Deff}{D_{\mathrm{bulk}}}
  = 1 + \frac{\dfrac{\fgb}{1-\fgb}\,( \alpha_{\mathrm{gb}} - 1)}
             {1 + \dfrac{2(\alpha_{\mathrm{gb}}-1)\,\fgb}{3(1-\fgb)}}.
  \label{eq:hart}
\end{equation}

\paragraph{Oriani trap occupancy.} Traps (dislocations, vacancies, GBs,
precipitates) immobilise hydrogen. Under local equilibrium
\citep{oriani1970}, each trap type $i$ has an equilibrium constant and
occupancy
\begin{equation}
  K_i = \exp\!\left(\frac{E_{b,i}}{RT}\right), \qquad
  \theta_{T,i} = \frac{K_i\,\theta_L}{1 + K_i\,\theta_L}, \qquad
  \theta_L = \frac{N_A C_L}{N_L},
  \label{eq:orianitheta}
\end{equation}
where $E_{b,i}$ is the trap binding energy [$\mathrm{J\,mol^{-1}}$]
(dislocations 20--30, GBs 40--50, vacancies 40--45, precipitates
60--90~$\mathrm{kJ\,mol^{-1}}$), $\theta_L$ the lattice-site occupancy [--],
$C_L$ the mobile (lattice) hydrogen concentration [$\mathrm{mol\,m^{-3}}$],
$N_L$ the interstitial site density [$\mathrm{m^{-3}}$]
($\approx 1.06\times10^{29}$ for FCC), and $N_A$ Avogadro's number. The model
is valid for $\theta_{T,i} < 0.9$; beyond that the kinetic McNabb--Foster
treatment \citep{mcnabb1963} is recommended.

\paragraph{Trapping-reduced diffusivity.} Partitioning hydrogen between mobile
and trapped populations gives the Oriani reduction
\begin{equation}
  \Deff = \frac{D}{1 + \displaystyle\sum_i \frac{N_{T,i}\,K_i}{N_L}},
  \label{eq:orianiDeff}
\end{equation}
where $N_{T,i}$ is the trap density of type $i$ [$\mathrm{m^{-3}}$]. The
denominator term $\sum_i N_{T,i}K_i/N_L$ is the ratio of trapped to mobile
hydrogen; the critical trap density at which $\Deff = D/2$ is
$N_T^{*} = N_L/K$.

\paragraph{Combined microstructure model.} Enhancement and trapping are
applied sequentially (GB first, then traps, assumed uniformly distributed):
\begin{equation}
  \Deff = \frac{(1-\fgb)\,D_{\mathrm{bulk}} + \fgb\,\alpha_{\mathrm{gb}} D_{\mathrm{bulk}}}
               {1 + \displaystyle\sum_i \frac{N_{T,i}\,K_i}{N_L}}.
  \label{eq:combined}
\end{equation}
The net modification factor $M = \Deff/D_{\mathrm{bulk}}$ can be either
$>1$ (GB-dominated, typically coarse traps / high $T$) or $<1$
(trapping-dominated, typically fine microstructure / low $T$)
\citep{oudriss2012}. Thermal vacancies, one of the trap populations, follow
\begin{equation}
  C_v = N \exp\!\left(-\frac{E_f^{v}}{RT}\right),
  \label{eq:vacancy}
\end{equation}
with $E_f^{v}$ the vacancy formation energy [$\mathrm{J\,mol^{-1}}$]
($140\;\mathrm{kJ\,mol^{-1}}$ for Incoloy~800) and $N$ the lattice site
density [$\mathrm{m^{-3}}$] \citep{kraftmakher1998}.

\paragraph{L4 permeation flux.} The membrane flux keeps the Sieverts boundary
conditions (solubility is a bulk thermodynamic property, unchanged by
microstructure) but replaces $D$ with $\Deff$:
\begin{equation}
  J_{\mathrm{L4}} = \frac{\Deff \Ks}{\Lm}\left(\sqrt{\Pup} - \sqrt{\Pdown}\right).
  \label{eq:L4flux}
\end{equation}
Because trap occupancy depends on the local concentration $C(x)$ --- the
profile linear in $\sqrt{P}$-space, written out as Eq.~\eqref{eq:profile} ---
$\Deff$ varies through the thickness; the code evaluates it on an $n$-point
profile and averages (arithmetic, harmonic, inlet, or outlet --- the
\code{method} argument) before applying Eq.~\eqref{eq:L4flux}.

% ----------------------------------------------------------------------------
\subsection{L5: full system (defective oxide + defective metal)}
% ----------------------------------------------------------------------------

L5 has no separate physics: it composes the L3 oxide description with the L4
metal description through the same flux-continuity closure as L2b, now with a
concentration-dependent $\Deff$:
\begin{equation}
  \frac{\Dox\Kox}{\Lox}\bigl(\Pup - \Pint\bigr)
  =
  \frac{\Deff(\Pint)\,\Ks}{\Lm}\bigl(\sqrt{\Pint} - \sqrt{\Pdown}\bigr).
  \label{eq:closureL5}
\end{equation}
Since $\Deff$ depends on the concentration profile, which depends on $\Pint$,
the solver iterates: guess $\Pint$, evaluate $\Deff$ on the resulting profile,
re-solve Eq.~\eqref{eq:closureL5} by Brent's method, and repeat until $\Pint$
converges (\code{interface\_solver.py} $\rightarrow$
\code{solve\_interface\_pressure\_defective\_metal}). The defective-oxide
variant applies this closure per parallel path of Eq.~\eqref{eq:L3parallel};
the assembly is orchestrated in \code{Application/Proposal.ipynb}. Regime
classification reuses Eq.~\eqref{eq:chi} with $\Rmet$ evaluated at $\Deff$.

% ----------------------------------------------------------------------------
\subsection{Shared temperature dependence}
% ----------------------------------------------------------------------------

Every transport property in both models ($\Dm$, $\Ks$, $\Dox$, $\Kox$,
$\kdiss$, $\Keq$) is evaluated at temperature with the same
\emph{reference-temperature} Arrhenius form:
\begin{equation}
  k(T) = k_{\mathrm{ref}} \exp\!\left[-\frac{E}{R}\left(\frac{1}{T} - \frac{1}{T_{\mathrm{ref}}}\right)\right],
  \label{eq:arrhenius}
\end{equation}
where $k_{\mathrm{ref}}$ is the measured value at the reference temperature
$T_{\mathrm{ref}}$ [K] and $E$ the relevant activation or solution energy
[$\mathrm{J\,mol^{-1}}$]: $E_D$ for diffusivities, $H_s$ (metal) or
$H_{\mathrm{sol,ox}}$ (oxide) for solubilities, $E_{\mathrm{diss}}$ for the
dissociation rate constant. This form avoids extrapolating pre-exponentials
far outside the fitted range. The permeability inherits
\begin{equation}
  Q_\Phi = E_D + H_s,
  \qquad \Phi(T) = \Phi_{\mathrm{ref}}
  \exp\!\left[-\frac{Q_\Phi}{R}\left(\frac1T - \frac1{T_{\mathrm{ref}}}\right)\right],
  \label{eq:QPhi}
\end{equation}
so an Arrhenius plot of measured permeation flux yields $Q_\Phi$ directly ---
the basis of the Arrhenius validation panels in both notebooks. Material
parameters ($D_{\mathrm{ref}}$, $E_D$, $K_{s,\mathrm{ref}}$, $H_s$, oxide
analogues, surface kinetics) live in \code{calculations/config/model\_config.py}
and \code{data/}; Section~\ref{sec:parameters} lists every value with its
default, sensitivity range, and literature source
\citep{schmidt1985,nemanic2023,stover1986,chen2011,grant1988,zhu2021,lu2022,young1997}.

% ============================================================================
\section{Model B: surface-kinetics extension (L6 and combinations)}
\label{sec:modelB}
% ============================================================================

Model B inserts a finite dissociation/recombination step at the gas--solid
interface, in series with the Model A bulk layers. All of Model B lives in
\code{calculations/surface\_kinetics.py}; the bulk equations
\eqref{eq:sieverts}--\eqref{eq:QPhi} are reused unchanged.

% ----------------------------------------------------------------------------
\subsection{Summary table}
% ----------------------------------------------------------------------------

{\footnotesize
\setlength{\tabcolsep}{4pt}
\begin{longtable}{@{}c >{\raggedright\arraybackslash}p{2.4cm} l >{\raggedright\arraybackslash}p{4.4cm} >{\raggedright\arraybackslash}p{2.2cm}@{}}
\caption{Model B equations, code locations, and references. All code locations
are functions in \code{calculations/surface\_kinetics.py} unless noted.}
\label{tab:summaryB}\\
\toprule
Eq. & Name & Formula & Code location & Ref. \\
\midrule
\endfirsthead
\toprule
Eq. & Name & Formula & Code location & Ref. \\
\midrule
\endhead
\bottomrule
\endfoot
\multicolumn{5}{@{}l}{\emph{L6 fundamentals}}\\
\eqref{eq:ads} & Dissociative adsorption & $r_{\mathrm{ads}} = \kdiss P (1-\theta)^2$ & \code{surface\_flux} & \citet{pick1985,baskes1980} \\
\eqref{eq:des} & Recombinative desorption & $r_{\mathrm{des}} = (\kdiss/\Keq)\,\theta^2$ & \code{surface\_flux} & \citet{pick1985,wampler1986} \\
\eqref{eq:Jsurf} & Net surface flux & $J_{\mathrm{surf}} = \kdiss\bigl[P(1-\theta)^2 - \theta^2/\Keq\bigr]$ & \code{surface\_flux} & \citet{andrew1992} \\
\eqref{eq:langmuir} & Langmuir isotherm & $\theta_{\mathrm{eq}} = \frac{\sqrt{\Keq P}}{1+\sqrt{\Keq P}}$ & implicit ($J_{\mathrm{surf}}=0$) & \citet{langmuir1918} \\
\eqref{eq:gtheta} & Coverage conversion & $g(\theta) = \frac{\theta}{(1-\theta)\sqrt{\Keq}}$ & \code{g\_theta} & \\
\eqref{eq:Rsurface} & Surface resistance & $\Rsurf \approx \frac{1}{\kdiss \Pup (1-\theta)^2}$ & \code{smooth\_surface\_}\newline\code{resistance} & \\
\eqref{eq:permeances} & Permeances & $\alpha = \frac{\Dox\Kox}{\Lox}$, $\beta = \frac{\Dm\Ks}{\Lm}$ & \code{compute\_permeances} & \\
\eqref{eq:KsKeq} & Sieverts--Langmuir link & $\Ks = \sqrt{\kdiss/\krec}\;N_{\mathrm{surf}}/N_L$ & \code{data/surface\_}\newline\code{kinetics\_data.py} & \citet{pick1985} \\
\midrule
\multicolumn{5}{@{}l}{\emph{L1+L6 --- metal + surface}}\\
\eqref{eq:JmL1L6} & Metal flux & $J_{\mathrm{m}} = \beta\bigl(g(\theta) - \sqrt{\Pdown}\bigr)$ & \code{metal\_flux\_L1} & \\
\eqref{eq:closureL1L6} & Closure in $\theta$ & $J_{\mathrm{surf}}(\theta) = J_{\mathrm{m}}(\theta)$ & \code{solve\_steady\_state\_}\newline\code{flux\_L1L6} & \\
\midrule
\multicolumn{5}{@{}l}{\emph{L2a+L6 --- oxide + surface}}\\
\eqref{eq:JoxL2aL6} & Oxide flux & $J_{\mathrm{ox}} = \alpha\bigl(g(\theta) - \sqrt{\Pdown}\bigr)$ & \code{oxide\_flux\_L2a} & \\
\midrule
\multicolumn{5}{@{}l}{\emph{L2+L6 --- oxide + metal + surface}}\\
\eqref{eq:PintTheta} & Analytic interface pressure & $\sqrt{\Pint} = \frac{\alpha g(\theta) + \beta\sqrt{\Pdown}}{\alpha+\beta}$ & \code{sqrt\_P\_int\_from\_theta} & \\
\eqref{eq:tripleflux} & Three-flux system & $J_{\mathrm{surf}} = J_{\mathrm{ox}} = J_{\mathrm{m}}$ & \code{solve\_steady\_state\_flux} & \\
\eqref{eq:fractions} & Resistance fractions & $f_i = R_i / R_{\mathrm{total}}$ & \code{solve\_steady\_state\_flux} & \\
\midrule
\multicolumn{5}{@{}l}{\emph{L3+L6 --- defective oxide + surface}}\\
\eqref{eq:alphapaths} & Path permeances & $\alpha_{\mathrm{crack}} = \alpha/\gamma$, $\alpha_{\mathrm{gb}} = \delta_D\alpha$, $\alpha_{\mathrm{pin}} = \infty$ & \code{calculate\_parallel\_}\newline\code{path\_flux\_L6} & \citet{strehlow1974} \\
\eqref{eq:L3L6total} & Multi-path total & $J_{\mathrm{tot}} = f_{\mathrm{int}}J_{\mathrm{int}} + \sum_i f_i J_i$ & \code{calculate\_mixed\_}\newline\code{defect\_flux\_L6} & \\
\midrule
\multicolumn{5}{@{}l}{\emph{L2+L4+L6 --- oxide + defective metal + surface}}\\
\eqref{eq:betaeff} & Effective permeance & $\beta_{\mathrm{eff}} = \Deff\Ks/\Lm$ & \code{calculate\_defective\_}\newline\code{metal\_flux\_L6} & \\
\eqref{eq:profile} & Concentration profile & $C(x) = \Ks\bigl[\sqrt{\Pint} - (\sqrt{\Pint}-\sqrt{\Pdown})\frac{x}{\Lm}\bigr]$ & \code{calculate\_defective\_}\newline\code{metal\_flux\_L6} & \\
\midrule
\multicolumn{5}{@{}l}{\emph{L3+L4+L6 --- full model}}\\
\eqref{eq:L346} & Full-model total & $J_{\mathrm{tot}} = f_{\mathrm{int}}J_{\mathrm{int}} + \sum_i f_i J_i$ & \code{calculate\_full\_model\_}\newline\code{flux\_L346\_v2} & \\
\end{longtable}
}

% ----------------------------------------------------------------------------
\subsection{L6 fundamentals: Langmuir--Hinshelwood surface kinetics}
% ----------------------------------------------------------------------------

\paragraph{Dissociative adsorption.} An $\mathrm{H_2}$ molecule needs two
adjacent vacant surface sites to dissociate,
$\mathrm{H_2(g)} + 2S \rightarrow 2\mathrm{H(ads)}$, giving a rate
second-order in the vacancy fraction:
\begin{equation}
  r_{\mathrm{ads}} = \kdiss\, P\, (1-\theta)^2,
  \label{eq:ads}
\end{equation}
where $\theta \in [0,1]$ is the surface coverage (fraction of occupied
sites) [--], $\kdiss$ the dissociation rate constant
[$\mathrm{mol\,m^{-2}\,s^{-1}\,Pa^{-1}}$], and $P$ the local $\mathrm{H_2}$
pressure [Pa] \citep{pick1985,baskes1980}. $\kdiss$ follows
Eq.~\eqref{eq:arrhenius} with activation energy $E_{\mathrm{diss}}$; the
numerical rate constants for both the oxide surface and the clean metal
surface (used for pinhole paths) are taken from the 316-steel measurements of
\citet{grant1988}, the only literature source providing the full dissociation
constant $k_1$ under clean, partially cleaned, and oxidised surface
conditions.

\paragraph{Recombinative desorption.} The reverse step,
$2\mathrm{H(ads)} \rightarrow \mathrm{H_2(g)} + 2S$, is second-order in
coverage. The code parameterises the recombination constant through the
surface equilibrium constant $\Keq = \kdiss/\krec$ [$\mathrm{Pa^{-1}}$]:
\begin{equation}
  r_{\mathrm{des}} = \krec\,\theta^2 = \frac{\kdiss}{\Keq}\,\theta^2 .
  \label{eq:des}
\end{equation}

\paragraph{Net surface flux.} The net rate of atoms entering the solid is the
Langmuir--Hinshelwood balance
\begin{equation}
  J_{\mathrm{surf}}(\theta)
  = \kdiss \left[\Pup\,(1-\theta)^2 - \frac{\theta^2}{\Keq}\right],
  \label{eq:Jsurf}
\end{equation}
positive for net adsorption \citep{andrew1992}. This is the extra series
element that Model B adds to every bulk configuration.

\paragraph{Langmuir isotherm.} Setting $J_{\mathrm{surf}} = 0$ recovers the
equilibrium coverage
\begin{equation}
  \frac{\theta_{\mathrm{eq}}}{1-\theta_{\mathrm{eq}}} = \sqrt{\Keq P}
  \quad\Longrightarrow\quad
  \theta_{\mathrm{eq}} = \frac{\sqrt{\Keq P}}{1 + \sqrt{\Keq P}},
  \label{eq:langmuir}
\end{equation}
which grows as $\sqrt{P}$ at low pressure and saturates at
$\theta_{\mathrm{eq}} \to 1$ \citep{langmuir1918}. When the surface is fast
($\theta \to \theta_{\mathrm{eq}}$) Model B collapses onto Model A --- this
limit is used as a verification test in
\code{Application/Surface\_proposal.ipynb}.

\paragraph{Coverage-to-driving-force conversion.} Inverting
Eq.~\eqref{eq:langmuir} defines the effective $\sqrt{P}$ that a coverage
$\theta$ presents to the bulk layers beneath:
\begin{equation}
  g(\theta) = \frac{\theta}{(1-\theta)\sqrt{\Keq}} \qquad [\mathrm{Pa^{1/2}}].
  \label{eq:gtheta}
\end{equation}
$g(\theta)$ is the pivot of the whole Model B formulation: it lets every bulk
flux be written as a function of the single unknown $\theta$, so the coupled
surface--bulk problem reduces to one-dimensional root finding. It also
connects to the Sieverts constant: consistency of the surface equilibrium
with bulk solubility requires
\begin{equation}
  \Ks = \sqrt{\kdiss/\krec}\;\frac{N_{\mathrm{surf}}}{N_L},
  \label{eq:KsKeq}
\end{equation}
with $N_{\mathrm{surf}} \approx 10^{19}\;\mathrm{m^{-2}}$ the surface site
density \citep{pick1985}.

\paragraph{Surface resistance and permeances.} Linearising
Eq.~\eqref{eq:Jsurf} gives the surface resistance used for regime
classification,
\begin{equation}
  \Rsurf \approx \frac{1}{\kdiss\,\Pup\,(1-\theta)^2}
  \qquad [\mathrm{Pa^{1/2}\,s\,m^{2}\,mol^{-1}}],
  \label{eq:Rsurface}
\end{equation}
which blows up as the surface saturates ($\theta \to 1$) or as $\kdiss\Pup$
becomes small --- the two routes to surface limitation. The bulk layers are
condensed into the permeances
\begin{equation}
  \alpha = \frac{\Dox\Kox}{\Lox},
  \qquad
  \beta = \frac{\Dm\Ks}{\Lm},
  \qquad
  \text{both } [\mathrm{mol\,m^{-2}\,s^{-1}\,Pa^{-1/2}}],
  \label{eq:permeances}
\end{equation}
with bulk resistances $\Rox = 1/\alpha$ and $\Rmet = 1/\beta$.

\paragraph{Change of oxide dissolution convention.}
\label{sec:oxideconvention}
Note a deliberate difference from Model A: there, dissociation happens (fast)
at the oxide--metal interface, so the oxide carries molecular $\mathrm{H_2}$
and obeys Henry's law (Eq.~\eqref{eq:henry}, driving force $\Delta P$). In
Model B, dissociation happens at the \emph{gas--oxide} surface
(Eq.~\eqref{eq:Jsurf}), so it is atomic hydrogen that enters the oxide and
all bulk driving forces are expressed in $\sqrt{P}$ units --- the code
docstrings accordingly call $\Kox$ the ``oxide Sieverts constant''
[$\mathrm{mol\,m^{-3}\,Pa^{-1/2}}$]. This makes $\alpha$ and $\beta$
dimensionally identical, which is what allows the permeance-weighted
interface pressure of Eq.~\eqref{eq:PintTheta}.

% ----------------------------------------------------------------------------
\subsection{L1+L6: perfect metal + surface kinetics}
% ----------------------------------------------------------------------------

With no oxide, the coverage sits directly on the metal, so the interface
driving force is $\sqrt{\Pint} = g(\theta)$ and the metal flux is
\begin{equation}
  J_{\mathrm{m}}(\theta) = \beta\bigl(g(\theta) - \sqrt{\Pdown}\bigr).
  \label{eq:JmL1L6}
\end{equation}
Steady state is the single-unknown closure
\begin{equation}
  J_{\mathrm{surf}}(\theta) = J_{\mathrm{m}}(\theta)
  \quad\Longrightarrow\quad \theta_{\mathrm{ss}},
  \qquad
  J_{\mathrm{ss}} = \beta\bigl(g(\theta_{\mathrm{ss}}) - \sqrt{\Pdown}\bigr),
  \label{eq:closureL1L6}
\end{equation}
solved by Brent's method on $\theta \in (0,1)$. The physics question this
level answers: at what $(\Pup, T)$ does control shift from dissociation
($\Rsurf > \Rmet$; flux $\propto \Pup$, slope 1.0) to bulk diffusion
($\Rmet > \Rsurf$; Sieverts recovered, slope 0.5)? High $\kdiss$ must
reproduce Eq.~\eqref{eq:L1flux} exactly --- the L1+L6 parity check in the
notebook.

% ----------------------------------------------------------------------------
\subsection{L2a+L6: perfect oxide + surface kinetics}
% ----------------------------------------------------------------------------

The same structure with the oxide as the only bulk layer:
\begin{equation}
  J_{\mathrm{ox}}(\theta) = \alpha\bigl(g(\theta) - \sqrt{\Pdown}\bigr),
  \qquad
  J_{\mathrm{surf}}(\theta) = J_{\mathrm{ox}}(\theta)
  \;\Longrightarrow\; \theta_{\mathrm{ss}}.
  \label{eq:JoxL2aL6}
\end{equation}
Resistances $\Rsurf$ (Eq.~\eqref{eq:Rsurface}) and $\Rox = 1/\alpha$ decide
whether the film or its surface is limiting.

% ----------------------------------------------------------------------------
\subsection{L2+L6: oxide + metal + surface kinetics}
% ----------------------------------------------------------------------------

\paragraph{Analytic reduction to one unknown.} Three fluxes must match:
surface (Eq.~\eqref{eq:Jsurf}), oxide, and metal. Nominally the unknowns are
$\theta$ and $\Pint$, but equating the oxide and metal fluxes,
$\alpha\bigl(g(\theta) - \sqrt{\Pint}\bigr) =
\beta\bigl(\sqrt{\Pint} - \sqrt{\Pdown}\bigr)$, is \emph{linear} in
$\sqrt{\Pint}$ and solves analytically:
\begin{equation}
  \sqrt{\Pint}(\theta)
  = \frac{\alpha\, g(\theta) + \beta \sqrt{\Pdown}}{\alpha + \beta}.
  \label{eq:PintTheta}
\end{equation}
The interface pressure is the permeance-weighted mean of the two boundary
driving forces. Substituting back leaves a single residual in $\theta$:
\begin{equation}
  J_{\mathrm{surf}}(\theta)
  = J_{\mathrm{ox}}(\theta)
  = J_{\mathrm{m}}(\theta), \qquad
  J_{\mathrm{ox}} = \alpha\bigl(g(\theta) - \sqrt{\Pint}(\theta)\bigr), \quad
  J_{\mathrm{m}} = \beta\bigl(\sqrt{\Pint}(\theta) - \sqrt{\Pdown}\bigr),
  \label{eq:tripleflux}
\end{equation}
solved with Brent's method on $f(\theta) = J_{\mathrm{surf}} -
J_{\mathrm{ox}} = 0$. This reduction (rather than nested 2-D solving) is what
keeps the higher combinations tractable.

\paragraph{Rate-limiting analysis.} Three series resistances,
$R_{\mathrm{total}} = \Rsurf + 1/\alpha + 1/\beta$, give fractional
contributions
\begin{equation}
  f_i = \frac{R_i}{R_{\mathrm{total}}}, \qquad
  i \in \{\text{surface},\, \text{oxide},\, \text{metal}\},
  \label{eq:fractions}
\end{equation}
and the step with $f_i > 0.5$ is declared rate-limiting (otherwise
\emph{mixed}). The pressure-exponent signatures --- surface-limited
$J \propto P^{1.0}$, oxide- or metal-limited $J \propto P^{0.5}$ under
Model B's $g(\theta)$ formulation --- are used in the notebook's regime maps
over $(\Pup, T, \kdiss)$.

% ----------------------------------------------------------------------------
\subsection{L3+L6: defective oxide + surface kinetics}
% ----------------------------------------------------------------------------

Each oxide path from L3 becomes an independent L2+L6 problem with its own
permeance,
\begin{equation}
  \alpha_{\mathrm{intact}} = \frac{\Dox\Kox}{\Lox}, \qquad
  \alpha_{\mathrm{crack}} = \frac{\alpha_{\mathrm{intact}}}{\gamma}, \qquad
  \alpha_{\mathrm{gb}} = \delta_D\,\alpha_{\mathrm{intact}}, \qquad
  \alpha_{\mathrm{pinhole}} = \infty,
  \label{eq:alphapaths}
\end{equation}
($\gamma$, $\delta_D$ as in Eqs.~\eqref{eq:crack}--\eqref{eq:gbdefect}). Each
path is solved for its own coverage $\theta_i$ and interface pressure
$\Pint{}_{,i}$ via Eqs.~\eqref{eq:PintTheta}--\eqref{eq:tripleflux}; the
pinhole path replaces the oxide kinetics by metal surface kinetics
($\kdiss^{\mathrm{m}}$, $\Keq^{\mathrm{m}}$, so $\sqrt{\Pint} = g_{\mathrm{m}}(\theta)$)
or, if none are supplied, by the Sieverts limit $\sqrt{\Pint} = \sqrt{\Pup}$.
The total is the area-weighted sum over intact plus all defect types,
\begin{equation}
  J_{\mathrm{total}} = f_{\mathrm{intact}}\,J_{\mathrm{intact}}
                     + \sum_i f_i\, J_i,
  \qquad f_{\mathrm{intact}} = 1 - \sum_i f_i,
  \label{eq:L3L6total}
\end{equation}
with $f_i$ the area fraction of defect type $i$ [--].

% ----------------------------------------------------------------------------
\subsection{L2+L4+L6: oxide + defective metal + surface kinetics}
% ----------------------------------------------------------------------------

The L2+L6 solver is wrapped in a self-consistency loop over the metal
diffusivity. The metal permeance becomes
\begin{equation}
  \beta_{\mathrm{eff}} = \frac{\Deff \Ks}{\Lm},
  \label{eq:betaeff}
\end{equation}
with $\Deff$ from the combined microstructure model
(Eq.~\eqref{eq:combined}). Because trap occupancy depends on the local
concentration, $\Deff$ is evaluated on the through-thickness profile
\begin{equation}
  C(x) = \Ks\left[\sqrt{\Pint} - \bigl(\sqrt{\Pint} - \sqrt{\Pdown}\bigr)\frac{x}{\Lm}\right],
  \qquad 0 \le x \le \Lm,
  \label{eq:profile}
\end{equation}
averaged (arithmetic/harmonic/inlet/outlet), and iterated: solve
Eq.~\eqref{eq:tripleflux} with the current $\beta_{\mathrm{eff}}$, rebuild the
profile, update $\Deff$, repeat until the relative change in $\Deff$ falls
below tolerance. Outputs include the modification factor
$M = \Deff/D_{\mathrm{lattice}}$ and the resistance fractions of
Eq.~\eqref{eq:fractions} evaluated with $\beta_{\mathrm{eff}}$.

% ----------------------------------------------------------------------------
\subsection{L3+L4+L6: full model}
% ----------------------------------------------------------------------------

The most complete configuration composes everything: each oxide path of
Eq.~\eqref{eq:alphapaths} is solved as an L2+L4+L6 problem (surface kinetics
$\rightarrow$ path-specific $\alpha$ $\rightarrow$ defective metal with
iterated $\Deff$), then area-weighted:
\begin{equation}
  J_{\mathrm{total}} = f_{\mathrm{intact}}\,J_{\mathrm{intact}}^{\mathrm{L2+L4+L6}}
                     + \sum_i f_i\, J_i^{\mathrm{L2+L4+L6}}.
  \label{eq:L346}
\end{equation}
Each path carries its own $\theta_i$, $\Pint{}_{,i}$, and $\Deff{}_{,i}$; the
system-level rate-limiting step is taken from the path that dominates the
total flux (or reported as \emph{mixed} if no path contributes more than
70\%).

% ============================================================================
\section{Parameter values and sources}
\label{sec:parameters}
% ============================================================================

Both models are driven by a single unified configuration,
\code{calculations/config/model\_config.py}: the metal is X40 NiCrAlTi 31/19
(Incoloy~802) with transport data from \citet{schmidt1985}, the oxide is
$\mathrm{Cr_2O_3}$ (``sample~4'', the undamaged-chromia limit) assembled from
\citet{nemanic2023,stover1986,chen2011}, surface kinetics are the 316-steel
measurements of \citet{grant1988}, and the metal microstructure (grain size,
trap spectrum) is parameterised for Hastelloy-N-class Ni alloys from
\citet{zhu2021,lu2022,young1997}. Tables~\ref{tab:paramsA}--\ref{tab:paramsC}
list every parameter with its default, the range explored in the sensitivity
analyses (\code{data/sensitivity\_parameters.csv}), the source, and why that
range was chosen. Entries marked \textsc{placeholder} have no literature
measurement; their defaults are assumptions and their ranges are deliberately
wide. Primary-source raw data are archived in the literature-database
workbook \code{permeation\_literature\_database\_v3.xlsx} in \code{data/}.

% ----------------------------------------------------------------------------
\subsection{Bulk transport and operating conditions}
% ----------------------------------------------------------------------------

{\footnotesize
\setlength{\tabcolsep}{4pt}
\begin{longtable}{@{}l l l >{\raggedright\arraybackslash}p{5.2cm}@{}}
\caption{Bulk transport parameters and operating conditions (used by both
models). Defaults from \code{model\_config.py}; ranges from the sensitivity
configuration.}
\label{tab:paramsA}\\
\toprule
Parameter & Default & Range & Source and range justification \\
\midrule
\endfirsthead
\toprule
Parameter & Default & Range & Source and range justification \\
\midrule
\endhead
\bottomrule
\endfoot
\multicolumn{4}{@{}l}{\emph{Metal (Incoloy~802) --- $\Dm$, $\Ks$ in Eqs.~\eqref{eq:sieverts}, \eqref{eq:L1flux}}}\\
$D_{\mathrm{ref}}$ \code{D\_ref} & $4.285\times10^{-9}\;\mathrm{m^2\,s^{-1}}$ & $[4.3\times10^{-11},\,4.3\times10^{-7}]$ & \citet{schmidt1985}: $\mathrm{T_2}$ permeation, X40 NiCrAlTi 31/19, 750--950\,$^\circ$C, 5\,mm plate. $\pm2$ decades covers alloy-to-alloy scatter across the Ni--Cr family. \\
$E_D$ \code{E\_D} & $72\,200\;\mathrm{J\,mol^{-1}}$ & $[60,\,80]\;\mathrm{kJ\,mol^{-1}}$ & \citet{schmidt1985}. Window brackets reported diffusion activation energies for austenitic Ni--Cr alloys. \\
$K_{s,\mathrm{ref}}$ \code{K\_s\_ref} & $0.05093\;\mathrm{mol\,m^{-3}\,Pa^{-1/2}}$ & $[10^{-4},\,10^{-1}]$ & \citet{schmidt1985}. Three-decade window spans reported Sieverts constants for austenitic alloys. \\
$H_s$ \code{H\_s} & $7\,200\;\mathrm{J\,mol^{-1}}$ & $[1,\,50]\;\mathrm{kJ\,mol^{-1}}$ & \citet{schmidt1985}. Solution enthalpy is small and weakly constrained; wide window for endothermic dissolvers. \\
$T_{\mathrm{ref}}$ (metal) & $1223\;\mathrm{K}$ & $[600,\,1400]\;\mathrm{K}$ & \citet{schmidt1985} measurement reference (950\,$^\circ$C). \\
\midrule
\multicolumn{4}{@{}l}{\emph{Oxide ($\mathrm{Cr_2O_3}$ sample~4) --- $\Dox$, $\Kox$ in Eqs.~\eqref{eq:henry}, \eqref{eq:L2aflux}}}\\
$D_{\mathrm{ox,ref}}$ \code{D\_ox\_ref} & $7.80\times10^{-19}\;\mathrm{m^2\,s^{-1}}$ & $[7.8\times10^{-21},\,7.8\times10^{-17}]$ & \citet{nemanic2023} ($\Phi_{\mathrm{ox}}$, $D$); \citet{stover1986} ($Q_p$, Fig.~7 sample Cr2); \citet{chen2011} ($E_D$). Undamaged-chromia limit; $\pm2$ decades for oxide quality/porosity variation. \\
$E_{D,\mathrm{ox}}$ \code{E\_D\_ox} & $70\,434\;\mathrm{J\,mol^{-1}}$ & $[69,\,71]\;\mathrm{kJ\,mol^{-1}}$ & \citet{chen2011} DFT (1.24\,eV), consistent with \citet{stover1986}. Tight range: theory and experiment agree. \\
$K_{\mathrm{ox,ref}}$ \code{K\_ox\_ref} & $0.35417\;\mathrm{mol\,m^{-3}\,Pa^{-1/2}}$ & $[0.08,\,0.4]$ & \citet{nemanic2023} + \citet{stover1986}: derived as $\Phi_{\mathrm{ox}}/\Dox$ at $T_{\mathrm{ref}}=673$\,K (Sieverts form, \S\ref{sec:oxideconvention}). \\
$H_{\mathrm{sol,ox}}$ \code{H\_sol\_ox} & $163\,566\;\mathrm{J\,mol^{-1}}$ & $[20,\,170]\;\mathrm{kJ\,mol^{-1}}$ & $Q_{p,\mathrm{ox}} - E_{D,\mathrm{ox}}$ from \citet{stover1986} and \citet{nemanic2023}; wide range because $Q_p$ measurements for chromia scatter strongly. \\
$T_{\mathrm{ref}}$ (oxide) & $673\;\mathrm{K}$ & $[600,\,1400]\;\mathrm{K}$ & \citet{nemanic2023} measurement reference (400\,$^\circ$C). \\
\midrule
\multicolumn{4}{@{}l}{\emph{Geometry and operating conditions --- $\Lm$, $\Lox$, $\Pup$, $T$}}\\
$\Lm$ \code{L\_metal} & $10^{-3}\;\mathrm{m}$ & $[5\times10^{-4},\,5\times10^{-3}]$ & Model default (\code{CONDITIONS}): 1\,mm plate; range spans practical component wall thicknesses. \\
$\Lox$ \code{thickness} & $4.8\times10^{-8}\;\mathrm{m}$ & $[10^{-10},\,10^{-5}]$ & Model default: sample-4 nominal 48\,nm scale; range spans native film to thick thermally grown scale. \\
$\Pup$ \code{P\_upstream} & $10^{5}\;\mathrm{Pa}$ & $[10^{-7},\,10^{7}]$ & Model default (\code{CONDITIONS}); range covers vacuum-facing to pressurised service. \\
$T$ \code{T\_operating} & $873\;\mathrm{K}$ & $[573,\,1273]\;\mathrm{K}$ & Model default (\code{CONDITIONS}, 600\,$^\circ$C); range spans the validated windows of the metal and oxide datasets. \\
\end{longtable}
}

% ----------------------------------------------------------------------------
\subsection{Oxide defects (L3) and metal microstructure (L4)}
% ----------------------------------------------------------------------------

{\footnotesize
\setlength{\tabcolsep}{4pt}
\begin{longtable}{@{}l l l >{\raggedright\arraybackslash}p{5.2cm}@{}}
\caption{Defect and microstructure parameters. Trap binding energies are
stored in \code{model\_config.py} as $\mathrm{eV}\times F$ with
$F = 96\,485\;\mathrm{C\,mol^{-1}}$.}
\label{tab:paramsB}\\
\toprule
Parameter & Default & Range & Source and range justification \\
\midrule
\endfirsthead
\toprule
Parameter & Default & Range & Source and range justification \\
\midrule
\endhead
\bottomrule
\endfoot
\multicolumn{4}{@{}l}{\emph{Oxide defects (L3) --- $f_i$, $\gamma$, $\delta_D$ in Eqs.~\eqref{eq:L3parallel}--\eqref{eq:mixed}}}\\
$f_{\mathrm{pinhole}}$ & $0.01$ & $[10^{-5},\,0.1]$ & Model default (\code{OXIDE\_DEFECTS}); range from trace porosity to heavily damaged scale. \\
$f_{\mathrm{crack}}$ & $0.005$ & $[10^{-5},\,0.1]$ & Model default (\code{OXIDE\_DEFECTS}). \\
$f_{\mathrm{gb,defect}}$ & $0.005$ & $[10^{-5},\,0.1]$ & Model default (\code{OXIDE\_DEFECTS}). \\
$\gamma$ \code{thickness\_factor} & $0.1$ & $[0.01,\,0.5]$ & Model default: crack retains 10\% of oxide thickness (Eq.~\eqref{eq:crack}); range from near-through-crack to half-thickness. \\
$\delta_D$ \code{diffusivity\_factor} & $10$ & $[1,\,100]$ & Model default: oxide-GB fast path (Eq.~\eqref{eq:gbdefect}); bounded by no enhancement and the fast-path limit. \\
\midrule
\multicolumn{4}{@{}l}{\emph{Metal microstructure (L4) --- $d$, $\delta$, $N_L$ in Eqs.~\eqref{eq:fgb}--\eqref{eq:combined}}}\\
$d$ \code{grain\_size} & $10^{-4}\;\mathrm{m}$ & $[10^{-5},\,5\times10^{-4}]$ & \citet{zhu2021}: Hastelloy-N equiaxed grains 50--150\,$\mu$m (EBSD); default is the midpoint, range extends to fine-grained condition. \\
$\delta$ \code{gb\_thickness} & $5\times10^{-10}\;\mathrm{m}$ & $[3\times10^{-10},\,10^{-9}]$ & Nominal high-angle boundary width; no Hastelloy-N measurement available. \\
$N_L$ \code{lattice\_density} & $8.774\times10^{28}\;\mathrm{m^{-3}}$ & $[8\times10^{28},\,1.2\times10^{29}]$ & Fe BCC estimate (model default); range spans BCC to FCC interstitial site densities. \\
\midrule
\multicolumn{4}{@{}l}{\emph{Trap spectrum (L4) --- $E_{b,i}$, $N_{T,i}$ in Eqs.~\eqref{eq:orianitheta}--\eqref{eq:orianiDeff}}}\\
$E_b$ dislocations & $19\,297\;\mathrm{J\,mol^{-1}}$ (0.20\,eV) & $[15,\,25]\;\mathrm{kJ\,mol^{-1}}$ & \citet{lu2022}: Alloy~625 TDS Peak~1, 0.186--0.215\,eV (dislocation elastic field); default is the mean. \\
$N_T$ dislocations & $8.16\times10^{12}\;\mathrm{m^{-3}}$ & $[8.2\times10^{10},\,8.2\times10^{14}]$ & \citet{zhu2021}: GND density from EBSD; $\pm2$ decades covers annealed to cold-worked states. \\
$E_b$ grain boundaries & $26\,051\;\mathrm{J\,mol^{-1}}$ (0.27\,eV) & $[25,\,30]\;\mathrm{kJ\,mol^{-1}}$ & \citet{lu2022}: TDS Peak~2, 0.258--0.281\,eV (GB + carbide interface). \\
$N_T$ grain boundaries & $6\times10^{14}\;\mathrm{m^{-3}}$ & $[6\times10^{12},\,6\times10^{16}]$ & Geometric estimate $3\rho_{\mathrm{gb}}/d$ for the 100\,$\mu$m equiaxed grain (Eq.~\eqref{eq:fgb}); $\pm2$ decades tracks the grain-size range. \\
$E_b$ vacancies & $41\,489\;\mathrm{J\,mol^{-1}}$ (0.43\,eV) & $[35,\,50]\;\mathrm{kJ\,mol^{-1}}$ & Ni estimate; no Hastelloy-N vacancy trap measurement. \\
$N_T$ vacancies & $10^{26}\;\mathrm{m^{-3}}$ & $[10^{24},\,10^{28}]$ & Model default; no measurement --- deliberately wide. \\
$E_b$ carbides ($\mathrm{M_6C}$) & $26\,051\;\mathrm{J\,mol^{-1}}$ (0.27\,eV) & $[20,\,40]\;\mathrm{kJ\,mol^{-1}}$ & \citet{lu2022}: M$_6$C TDS, 0.258--0.281\,eV; upper range allows for incoherent interfaces. \\
$N_T$ carbides & $2\times10^{25}\;\mathrm{m^{-3}}$ & $[2\times10^{23},\,2\times10^{27}]$ & \citet{young1997}: carbide trap density upper bound in Ni-17Cr-8Fe (TDS). \\
\end{longtable}
}

% ----------------------------------------------------------------------------
\subsection{Surface kinetics (L6)}
% ----------------------------------------------------------------------------

{\footnotesize
\setlength{\tabcolsep}{4pt}
\begin{longtable}{@{}l l l >{\raggedright\arraybackslash}p{4.1cm}@{}}
\caption{Surface-kinetics parameters (Model B). Both parameter sets derive
from \citet{grant1988}, the only literature source reporting the full
dissociation constant $k_1$; the oxidised-foil condition proxies the
$\mathrm{Cr_2O_3}$ surface and the clean condition the bare metal (pinhole
paths).}
\label{tab:paramsC}\\
\toprule
Parameter & Default & Range & Source and range justification \\
\midrule
\endfirsthead
\toprule
Parameter & Default & Range & Source and range justification \\
\midrule
\endhead
\bottomrule
\endfoot
\multicolumn{4}{@{}l}{\emph{Oxide surface (gas--$\mathrm{Cr_2O_3}$) --- $\kdiss$, $\Keq$ in Eqs.~\eqref{eq:ads}--\eqref{eq:Jsurf}}}\\
$k_{\mathrm{diss,ref}}$ \code{k\_diss\_ref} & $9.487\times10^{-8}\;\mathrm{mol\,m^{-2}\,s^{-1}\,Pa^{-1}}$ & $[9.5\times10^{-10},\,9.5\times10^{-6}]$ & \citet{grant1988}: oxidised 316 foil, $k_1 = 9.2\times10^{-8}$ at 965\,K, re-referenced to $T_{\mathrm{ref}}=1623$\,K; $\pm2$ decades because 316 oxide is a proxy for $\mathrm{Cr_2O_3}$. \\
$E_{\mathrm{diss}}$ \code{E\_diss} & $57\,950\;\mathrm{J\,mol^{-1}}$ & $[40,\,70]\;\mathrm{kJ\,mol^{-1}}$ & \citet{grant1988}: oxidised-both-surfaces condition ($E_{k_1} = 59\,861\;\mathrm{J\,mol^{-1}}$; default shifted slightly by the re-referencing). \\
$K_{\mathrm{eq,ref}}$ \code{K\_eq\_ref} & $10^{-4}\;\mathrm{Pa^{-1}}$ & $[10^{-14},\,10^{-2}]$ & \textsc{placeholder}: no gas--$\mathrm{Cr_2O_3}$ surface equilibrium constant in the literature; deliberately wide window bracketing the working default ($\pm2$ decades above, extending down to the historical sweep floor). \\
$H_{\mathrm{eq}}$ \code{H\_eq} & $20\,000\;\mathrm{J\,mol^{-1}}$ & $[5,\,50]\;\mathrm{kJ\,mol^{-1}}$ & \textsc{placeholder}: no measured surface equilibrium enthalpy for $\mathrm{Cr_2O_3}$. \\
$T_{\mathrm{ref}}$ (oxide surface) & $1623\;\mathrm{K}$ & $[600,\,1800]\;\mathrm{K}$ & Configuration choice: matches the $\mathrm{Cr_2O_3}$ transport reference temperature. \\
\midrule
\multicolumn{4}{@{}l}{\emph{Metal surface (pinhole paths, L1+L6) --- $\kdiss^{\mathrm{m}}$, $\Keq^{\mathrm{m}}$}}\\
$k_{\mathrm{diss,ref}}^{\mathrm{m}}$ \code{k\_diss\_metal\_ref} & $1.346\times10^{-6}\;\mathrm{mol\,m^{-2}\,s^{-1}\,Pa^{-1}}$ & $[1.3\times10^{-8},\,1.3\times10^{-4}]$ & \citet{grant1988}: clean 316 surface, $k_1$ at 965\,K; $\pm2$ decades for surface-condition uncertainty. \\
$E_{\mathrm{diss}}^{\mathrm{m}}$ \code{E\_diss\_metal} & $81\,560\;\mathrm{J\,mol^{-1}}$ & $[70,\,100]\;\mathrm{kJ\,mol^{-1}}$ & \citet{grant1988}: clean-surface activation energy. \\
$K_{\mathrm{eq,ref}}^{\mathrm{m}}$ \code{K\_eq\_metal\_ref} & $10^{-3}\;\mathrm{Pa^{-1}}$ & $[10^{-10},\,10^{-3}]$ & \textsc{placeholder}: no direct surface equilibrium measurement for Incoloy~802. \\
$H_{\mathrm{eq}}^{\mathrm{m}}$ \code{H\_eq\_metal} & $15\,000\;\mathrm{J\,mol^{-1}}$ & $[5,\,40]\;\mathrm{kJ\,mol^{-1}}$ & \textsc{placeholder}: no direct measurement. \\
$T_{\mathrm{ref}}$ (metal surface) & $965\;\mathrm{K}$ & $[600,\,1100]\;\mathrm{K}$ & \citet{grant1988} experimental reference temperature (692\,$^\circ$C). \\
\end{longtable}
}

% ============================================================================
\section{Cross-model equation map}
% ============================================================================

\begin{table}[h]
\centering
\caption{Which equations are shared between the two models and which are
specific. Shared equations are implemented once and imported by Model B.}
\label{tab:crossmodel}
\begin{tabular}{lcc>{\raggedright\arraybackslash}p{4.7cm}}
\toprule
Equation & Model A & Model B & Role in Model B \\
\midrule
Sieverts' law \eqref{eq:sieverts} & \checkmark & \checkmark & metal boundary condition via $g(\theta)$ \\
Fick's law \eqref{eq:fick} & \checkmark & \checkmark & unchanged bulk transport \\
Henry's law \eqref{eq:henry} & \checkmark & replaced & oxide carries atomic H --- Sieverts-form $\Kox$, $\sqrt{P}$ driving force (\S\ref{sec:oxideconvention}) \\
Arrhenius form \eqref{eq:arrhenius} & \checkmark & \checkmark & also applied to $\kdiss$, $\Keq$ \\
Parallel paths \eqref{eq:L3parallel} & \checkmark & \checkmark & per-path L2+L6 solves \eqref{eq:L3L6total} \\
Oriani trapping \eqref{eq:orianiDeff} & \checkmark & \checkmark & inside $\beta_{\mathrm{eff}}$ iteration \eqref{eq:betaeff} \\
GB enhancement \eqref{eq:gbparallel} & \checkmark & \checkmark & inside $\beta_{\mathrm{eff}}$ iteration \\
Interface closure \eqref{eq:closureL2b} & \checkmark & replaced & analytic $\sqrt{\Pint}(\theta)$ \eqref{eq:PintTheta} \\
Regime ratio $\chi$ \eqref{eq:chi} & \checkmark & generalised & three-way fractions \eqref{eq:fractions} \\
Surface kinetics \eqref{eq:ads}--\eqref{eq:Rsurface} & --- & \checkmark & the new series element \\
$g(\theta)$ conversion \eqref{eq:gtheta} & --- & \checkmark & reduces system to 1-D root find \\
\bottomrule
\end{tabular}
\end{table}

The essential difference is the closure variable: Model A solves for the
interface pressure $\Pint$ (Brent's method on
Eq.~\eqref{eq:closureL2b}/\eqref{eq:closureL5}), while Model B solves for the
surface coverage $\theta$ (Brent's method on Eq.~\eqref{eq:tripleflux}), with
$\Pint$ recovered analytically from Eq.~\eqref{eq:PintTheta}. In the
fast-surface limit ($\kdiss \to \infty$, $\theta \to \theta_{\mathrm{eq}}$)
every Model B level reduces to its Model A counterpart, which is the
consistency test used throughout \code{Application/Surface\_proposal.ipynb}.

% ============================================================================
\bibliographystyle{plainnat}
\bibliography{references}

% ============================================================================
\appendix
\section{Resistance formulation: complete derivation}
\label{app:resistance}
% ============================================================================

\begin{sloppypar}
This appendix derives the resistance quantities used in \S\ref{sec:modelA}
from first principles and explains precisely in what sense the metal and
oxide resistances can --- and cannot --- be compared. The corresponding code
lives in \code{oxide\_permeation.py} (functions
\code{calculate\_oxide\_resistance}, \code{calculate\_metal\_resistance},
and \code{calculate\_transition\_pressure}) and in
\code{interface\_solver.py} (\code{calculate\_oxide\_metal\_system}).
\end{sloppypar}

% ----------------------------------------------------------------------------
\subsection{What a permeation resistance is}
% ----------------------------------------------------------------------------

The electrical analogy defines a resistance through a linear response law,
\begin{equation}
  J = \frac{X_{\mathrm{up}} - X_{\mathrm{down}}}{R_X},
  \label{eq:app-linresp}
\end{equation}
where $X$ is some \emph{driving-force variable} evaluated at the two
boundaries, in analogy with $I = V/R$. The definition is only useful when $J$
is genuinely linear in $\Delta X$, because only then is $R_X$ a constant of
the layer, independent of operating point --- exactly as an ohmic resistor
has an $R$ independent of the applied voltage. The subtlety in this model is
that the two layers are linear in \emph{different} variables: the choice of
$X$ is not free, it is dictated by each layer's solubility law.

% ----------------------------------------------------------------------------
\subsection{Oxide: linear in $\Delta P$}
% ----------------------------------------------------------------------------

Combining Henry's law (Eq.~\eqref{eq:henry}) at both faces with Fick's first
law across the thickness:
\begin{equation}
  J_{\mathrm{ox}}
  = \Dox\,\frac{C_{\mathrm{up}} - C_{\mathrm{down}}}{\Lox}
  = \Dox\,\frac{\Kox \Pup - \Kox \Pdown}{\Lox}
  = \underbrace{\frac{\Dox\Kox}{\Lox}}_{1/\Rox}\,\bigl(\Pup - \Pdown\bigr).
  \label{eq:app-oxide}
\end{equation}
The flux is exactly linear in $\Delta P$, so the oxide possesses a true,
pressure-independent resistance
\begin{equation}
  \Rox = \frac{\Lox}{\Dox\Kox}
  \qquad [\mathrm{Pa\,s\,m^{2}\,mol^{-1}}],
  \label{eq:app-Rox}
\end{equation}
with the pressure difference itself as the driving force
($X = P$). Every extra pascal of driving force buys the same extra flux,
at any operating pressure.

% ----------------------------------------------------------------------------
\subsection{Metal: linear in $\Delta\sqrt{P}$}
% ----------------------------------------------------------------------------

The metal transports a different species. In the L2b picture the intact
$\mathrm{H_2}$ molecule arriving through the oxide dissociates at the
oxide--metal interface,
\begin{equation}
  \mathrm{H_2} \;\rightleftharpoons\; 2\,\mathrm{H}^{*},
  \label{eq:app-dissoc}
\end{equation}
and it is the \emph{atoms} that dissolve into and diffuse through the metal
lattice. The law of mass action for this equilibrium, with two atoms produced
per molecule, reads $C_{\mathrm{H}}^{2}/\Pint = \text{const}$, so the
dissolved-atom concentration scales as $C_{\mathrm{H}} = \Ks\sqrt{\Pint}$ ---
Sieverts' law (Eq.~\eqref{eq:sieverts}) is nothing but the 2:1 stoichiometry
of reaction~\eqref{eq:app-dissoc}. Repeating the resistance construction with
this boundary condition:
\begin{equation}
  J_{\mathrm{m}}
  = \Dm\,\frac{\Ks\sqrt{\Pup} - \Ks\sqrt{\Pdown}}{\Lm}
  = \underbrace{\frac{\Dm\Ks}{\Lm}}_{1/R_{\mathrm{L1}}}\,
    \bigl(\sqrt{\Pup} - \sqrt{\Pdown}\bigr).
  \label{eq:app-metal}
\end{equation}
The metal is exactly linear too --- but in $\Delta\sqrt{P}$, not $\Delta P$.
Its natural resistance is
\begin{equation}
  R_{\mathrm{L1}} = \frac{\Lm}{\Dm\Ks} = \frac{\Lm}{\Phi}
  \qquad [\mathrm{Pa^{1/2}\,s\,m^{2}\,mol^{-1}}],
  \label{eq:app-RL1}
\end{equation}
a constant of the layer, with driving force $X = \sqrt{P}$.

% ----------------------------------------------------------------------------
\subsection{Why the two are not directly comparable}
% ----------------------------------------------------------------------------

Both $\Rox$ and $R_{\mathrm{L1}}$ are legitimate constants, but they act on
different driving forces and carry different units:

\begin{center}
\begin{tabular}{lllll}
\toprule
Layer & Flux law & Driving force $X$ & Resistance & Units \\
\midrule
Oxide & $J = \Delta P/\Rox$ & $P$ & $\Rox = \Lox/(\Dox\Kox)$ & $\mathrm{Pa\,s\,m^2\,mol^{-1}}$ \\
Metal & $J = \Delta\sqrt{P}/R_{\mathrm{L1}}$ & $\sqrt{P}$ & $R_{\mathrm{L1}} = \Lm/(\Dm\Ks)$ & $\mathrm{Pa^{1/2}\,s\,m^2\,mol^{-1}}$ \\
\bottomrule
\end{tabular}
\end{center}

Asking whether $\Rox > R_{\mathrm{L1}}$ is dimensionally meaningless --- like
comparing a quantity in ohms with one in kilograms. Chemically, the mismatch
is the dissociation step \eqref{eq:app-dissoc} sitting between the two
layers: the oxide counts \emph{molecules} (one dissolved per gas molecule,
hence $\propto P$), the metal counts \emph{atoms} produced by an equilibrium
that delivers two per molecule (hence $\propto \sqrt{P}$). The unit
$\mathrm{Pa^{1/2}}$ separating the two resistances is exactly the
stoichiometric square root introduced by that reaction. Yet the L2b regime
question (``which layer limits the flux?'') demands a comparison. The
resolution is to re-express the metal on the oxide's per-pascal basis.

% ----------------------------------------------------------------------------
\subsection{Linearising the metal: origin of the factor 2}
% ----------------------------------------------------------------------------

Fix the downstream pressure and consider the metal flux as a function of its
upstream (interface) pressure, $J_{\mathrm{m}}(P) =
\tfrac{\Dm\Ks}{\Lm}\bigl(\sqrt{P} - \sqrt{\Pdown}\bigr)$. Taylor-expanding
about the operating point $P = \Pint$:
\begin{equation}
  J_{\mathrm{m}}(P) \approx J_{\mathrm{m}}(\Pint)
  + \frac{\mathrm{d}J_{\mathrm{m}}}{\mathrm{d}P}\bigg|_{\Pint}\,(P - \Pint),
  \qquad
  \frac{\mathrm{d}J_{\mathrm{m}}}{\mathrm{d}P}\bigg|_{\Pint}
  = \frac{\Dm\Ks}{\Lm}\cdot
    \frac{\mathrm{d}\sqrt{P}}{\mathrm{d}P}\bigg|_{\Pint}
  = \frac{\Dm\Ks}{\Lm}\cdot\frac{1}{2\sqrt{\Pint}}.
  \label{eq:app-taylor}
\end{equation}
The local slope $\mathrm{d}J/\mathrm{d}P$ is a per-pascal conductance; its
reciprocal is the metal's \emph{effective} per-pascal resistance at that
operating point,
\begin{equation}
  \Rmet(\Pint)
  = \left(\frac{\mathrm{d}J_{\mathrm{m}}}{\mathrm{d}P}\bigg|_{\Pint}\right)^{-1}
  = \frac{2\,\Lm\sqrt{\Pint}}{\Dm\Ks}
  \qquad [\mathrm{Pa\,s\,m^{2}\,mol^{-1}}],
  \label{eq:app-Rmet}
\end{equation}
which is Eq.~\eqref{eq:Rmetal}. The factor 2 is nothing but the derivative
$\mathrm{d}\sqrt{P}/\mathrm{d}P = 1/(2\sqrt{P})$ --- no additional physics.
The approximation is valid when the pressure drop across the metal is small
compared with $\Pint$, i.e.\ near the operating point about which the
expansion was taken. Physically: the dissociation equilibrium
\eqref{eq:app-dissoc} yields dissolved atoms only as $\sqrt{P}$, so at higher
pressure each additional pascal of $\mathrm{H_2}$ converts into
proportionally fewer new $\mathrm{H}^{*}$ atoms --- the marginal atom yield
falls as $1/(2\sqrt{\Pint})$. Each additional pascal therefore buys less
additional flux, and the metal's per-pascal resistance \emph{grows} as
$\sqrt{\Pint}$ while the oxide's, which counts molecules one-for-one, stays
constant.

% ----------------------------------------------------------------------------
\subsection{Consequences: regime ratio and transition pressure}
% ----------------------------------------------------------------------------

With both layers on the per-pascal basis, the ratio
$\chi = \Rox/\Rmet$ of Eq.~\eqref{eq:chi} is dimensionless and meaningful.
Since $\Rmet \propto \sqrt{\Pint}$, $\chi$ \emph{falls} as pressure rises:
the bilayer is oxide-limited at low pressure (flux slope $\to 1.0$) and
metal-limited at high pressure (slope $\to 0.5$). The crossover follows from
setting $\Rox = \Rmet$:
\begin{equation}
  \frac{\Lox}{\Dox\Kox} = \frac{2\,\Lm\sqrt{P_{\mathrm{trans}}}}{\Dm\Ks}
  \;\;\Longrightarrow\;\;
  \sqrt{P_{\mathrm{trans}}} = \frac{\Lox\,\Dm\Ks}{2\,\Dox\Kox\,\Lm}
  \;\;\Longrightarrow\;\;
  P_{\mathrm{trans}} = \left(\frac{\Lox\,\Dm\Ks}{2\,\Dox\Kox\,\Lm}\right)^{2},
  \label{eq:app-Ptrans}
\end{equation}
which is Eq.~\eqref{eq:Ptrans}.

% ----------------------------------------------------------------------------
\subsection{Caveats and the Model B alternative}
% ----------------------------------------------------------------------------

\paragraph{Resistances are diagnostics, not the solution.} Because the
series combination of a linear element ($\Rox$) and a linearised one
($\Rmet$) is only approximate, the L2b flux is \emph{never} computed as
$\Delta P/(\Rox + \Rmet)$. The code always solves the exact flux-continuity
condition Eq.~\eqref{eq:closureL2b} for $\Pint$ by Brent's method and only
afterwards evaluates $\Rox$, $\Rmet(\Pint)$, and $\chi$ to label the regime
(\code{interface\_solver.py} $\rightarrow$
\code{calculate\_oxide\_metal\_system}).

\paragraph{Model B sidesteps the mismatch.} In the surface-kinetics model
every layer is expressed in $\sqrt{P}$ units from the outset: dissociation at
the gas--oxide surface means atomic hydrogen traverses the oxide
(\S\ref{sec:oxideconvention}), so the permeances $\alpha$ and $\beta$ of
Eq.~\eqref{eq:permeances} are dimensionally identical and the resistances
$\Rsurf$, $1/\alpha$, $1/\beta$ add and compare directly
(Eq.~\eqref{eq:fractions}) --- no linearisation, and no factor of 2, is ever
needed there.


\end{document}
